% Source: https://artofproblemsolving.com/community/c3h2586378_advanced_linear_equations_prealgebra
% Stable identity: aops-topic-2586378-q02
The equation is equivalent to $(a-c)x=d-b$. If $a\ne c$, it has exactly one solution $x=\dfrac{d-b}{a-c}$. If $a=c$, it reads $0=d-b$, which is false for every $x$ when $b\ne d$, and true for every $x$ when $b=d$. So there are no solutions exactly when $\boxed{a=c\text{ and }b\ne d}$.
